\section{读出侧设计}\label{sec:readout}

第~\ref{sec:noise}~节说明，换平台、加功率都降不了理想探测器下的误差下限，能动的只有读出一侧和光源。
本节给出读出一侧的两种改法，看它们在现实噪声下能把误差降到多少、对激光的要求有什么影响，再讨论读出层的记忆从哪里来。

\subsection{加大下载端口耦合与放慢符号}

散粒噪声那一半下限取决于每个读数积分到的光子数。在不改变腔内动力学的前提下，有两种办法增加这个数：
\begin{itemize}
\item \emph{加大下载端口耦合}：把 $\Qtwo$ 的取值范围从固定的 $2.2\times10^{7}$ 放宽到 $10^{6}$ 到 $2.2\times10^{7}$，
      让腔内更大比例的光从下载端口出来到达探测器；代价是总品质因数降低，产生同样非线性所需的功率升高。
\item \emph{放慢符号}：把每个符号从 $1~\mathrm{ns}$ 延长到 $2~\mathrm{ns}$，每个读数积分的时间加倍，同时腔有更多时间跟上输入的变化；代价是速率减半。
\end{itemize}
两种改法都按第~\ref{sec:protocol}~节的流程在 256 个随机工作点中重新选点。
理想探测器下，它们把 Mackey--Glass 的误差下限从 $1.3\times10^{-4}$ 分别降到 $8.7\times10^{-5}$ 和 $6.9\times10^{-5}$；
现实噪声（PIN 探测器）下分别为 $1.55\times10^{-4}$ 和 $1.58\times10^{-4}$，比表~\ref{tab:encodings} 中的掩模（$3.3\times10^{-4}$）好约 2 倍。
加大下载端口耦合选出的点为 $Q=1.9\times10^{6}$、$178~\mathrm{mW}$（表~\ref{tab:encodings} 中用 PIN 探测器选出的普通掩模点为 $Q=6\times10^{6}$、$19~\mathrm{mW}$；第~\ref{sec:breakdown}~节用理想探测器选出的点为 $23~\mathrm{mW}$）；放慢符号选出的点为 $Q=5\times10^{6}$。
我们把加大下载端口耦合的点称为\emph{最好的单环}，以 $1.55\times10^{-4}$ 作为它的参考误差；第~\ref{sec:reference}、\ref{sec:laser}、\ref{sec:quant}~节用的都是它，
用其他噪声实现评估时为 $1.5$ 到 $1.6\times10^{-4}$。
与同读出维度的 ESN 相比，这两个结果在 ESN 输入噪声取 $3\times10^{-3}$ 时与它持平或略好（ESN 为 $1.9\times10^{-4}$ 到 $3.2\times10^{-4}$，取决于随机储备池的选择）；
输入噪声低于约 $2.5\times10^{-3}$ 时 ESN 更好（第~\ref{sec:software}~节）。

\subsection{对失谐抖动的敏感性}

两种改法降低的是散粒噪声那一半，另一半（失谐抖动）不受影响，所以改进后的误差对抖动更敏感。
图~\ref{fig:noise}b 把假设的抖动从 0 扫到半线宽的 3\%，每一档都重新选点。
加大下载端口耦合时，抖动为 0、0.1\%、0.3\%、1\%、3\% 的误差为
$5.7\times10^{-5}$、$6.4\times10^{-5}$、$8.1\times10^{-5}$、$1.6\times10^{-4}$、$2.8\times10^{-4}$；
放慢符号时为 $9.5\times10^{-5}$、$9.8\times10^{-5}$、$1.1\times10^{-4}$、$1.6\times10^{-4}$、$3.2\times10^{-4}$；
ESN（输入噪声 $3\times10^{-3}$）始终为 $3.2\times10^{-4}$。
抖动从 1\% 压到 0.3\% 和 0.1\%，加大下载端口耦合的误差分别降 1.9 和 2.4 倍；抖动到 3\% 时两种改法都退回到与 ESN 持平。

按式~\eqref{eq:linewidth}，加大下载端口耦合的工作点（$Q=1.9\times10^{6}$、$1~\mathrm{ns}$）上，0.1\%、0.3\%、1\%、3\% 的抖动分别对应约
$16~\mathrm{Hz}$、$150~\mathrm{Hz}$、$1.6~\mathrm{kHz}$、$14~\mathrm{kHz}$ 的激光本征线宽；
放慢符号的工作点 $Q$ 更高，同样的抖动对应的线宽约为前者的 0.3 倍（1\% 约 $470~\mathrm{Hz}$）。
两种改法对激光的要求不同。加大下载端口耦合是靠更低的 $Q$ 和更高的功率换来光子数，所以对同一台激光，它的相对抖动比普通掩模的点（$Q=6\times10^{6}$，1\% 只对应约 $160~\mathrm{Hz}$）小；
放慢符号靠更高的 $Q$ 和更长的符号，线宽要求比加大下载端口耦合的点约严 3 倍。在线宽为几十千赫兹以上的激光下，两者的好处都会被抖动吃掉（表~\ref{tab:laser}）。

\subsection{记忆从哪里来}\label{sec:memory}

微环自身的记忆只有光子寿命 $Q/(2\pi\nu_0)$（本文选出的工作点上为 $1.6$ 到 $5~\mathrm{ns}$，即几个符号），NARMA-10 却需要约 10 个符号的记忆。
我们在现实噪声下、用 31 根梳齿 $\times$ 12 个抽头的读出，分别试了增加微环自身的记忆和从外部加记忆两类办法：
\begin{itemize}
\item 增加微环自身的记忆没有用。不改时为 $0.18$（这组运行；表~\ref{tab:encodings} 中为 $0.19$）；把符号缩短到 $0.5$ 和 $0.25~\mathrm{ns}$，让光子寿命相对符号变长，为 $0.22$ 和 $0.20$；
      加一条把下载端口的输出经延迟一个符号的光纤环送回输入端的光反馈环路，为 $0.19$。
\item 从外部加记忆有用。去掉读出抽头时为 $0.42$，加上 12 个抽头降到 $0.19$；
      再让每根输入梳齿同时携带延迟 1 到 12 个符号的信号（输入延迟线，延迟长度随机抽取后按验证误差选，与 ESN 的选法相同），降到 $0.083$。
      但 ESN 加上同样的延迟线后为 $0.029$，差距没有缩小。
\end{itemize}
把几个掩模和延迟各不相同的环拼接读出后（训练长度加到 4 倍，仅此一项就让带延迟线的单环从 $0.083$ 降到 $0.037$；3 条序列 $\times$ 6 对新运行），
1、2、4、8 个环为 $0.037$、$0.013$、$0.011$、$0.012$，4 个环与带同样延迟线的 100 节点 ESN（$0.010$）持平；
不带延迟线的 100 和 300 节点 ESN 在这个训练长度下为 $0.015$ 和 $0.012$（在第~\ref{sec:software}~节的标准长度下为 $0.041$ 和 $0.032$）。
在需要长记忆的混沌预测任务 Lorenz（向前 10 步）上，同样的做法从单环的 $1.3\times10^{-2}$ 降到 8 个环的 $2.6\times10^{-3}$，
但仍比带延迟线的 100 节点 ESN（$1.05\times10^{-3}$）差约 2.5 倍。

这些结果提示第三条设计规则：\emph{微环贡献的主要是非线性特征而不是记忆；需要记忆的任务应当在外部加延迟线或读出抽头，再用多个掩模不同的环增加特征的多样性}。
用线性记忆容量和信息处理容量~\cite{dambre2012}直接测量，才能把这一点定量化，留作后续工作。
这与掩模在 Mackey--Glass（需要强非线性、短记忆）上优势最大、在 NARMA-10 上不成立是一致的。
（单环的这几组实验与表~\ref{tab:encodings} 用同一流程：每种变体 256 个随机工作点、3 条序列、6 对新运行复测；
读出抽头的那组对比 $0.42$ 与 $0.19$ 取自同一次搜索中选点时的数值。）

\subsection{非线性读出与更长的记忆}\label{sec:nonlin-readout}

前面各节的读出对取对数后的梳齿功率是线性的。加上当前 31 路特征两两的乘积（多 496 个特征），并用 4 次独立噪声的运行（而不是 1 次）拟合，
最好单环的误差从 $1.7\times10^{-4}$ 降到 $2.4\times10^{-5}$，好 7 倍，无抖动误差从 $5.3\times10^{-5}$ 降到 $8.9\times10^{-6}$；抖动损失仍约为 3。
同样的读出只让 31 节点 ESN（输入噪声 $3\times10^{-3}$）从 $2.7\times10^{-4}$ 降到 $1.8\times10^{-4}$，放到用同样二次读出的 31 节点 ESN 的误差--输入噪声曲线上，微环从约 $1.7\times10^{-3}$ 的输入噪声移到约 $2.3\times10^{-4}$，按第~\ref{sec:software}~节的换算约多 3 位。
31 路读数量化后，这一优势大部分还在（对数前端；ESN 量化后的输入加上抖动副本，使它也看到多次噪声实现）：6、7 位时微环比 ESN 好 1.5 到 1.6 倍，8 到 10 位持平，12 位时差 2.2 倍。

读出的记忆加长后，这一优势就消失了。把抽头从 8 个增加到 16 或 32 个，并加入当前与上一符号特征之间的乘积，两者都会改善，但 ESN 改善更多：
8 位时持平（32 个抽头、用 8 次运行拟合时微环为 $8.9\times10^{-6}$，与上文的无抖动误差恰好相同；ESN 为 $7.1\times10^{-6}$），10 位时 ESN 好 3 到 4 倍（32 个抽头时 $8.2\times10^{-7}$ 对 $3.5\times10^{-6}$）。
所以更丰富的读出能把微环的误差降一个数量级以上，但不改变它相对软件的位置；为使结果可比，其余各处仍用 8 抽头的线性读出。
